\chapter{Poisson Brackets, Canonical Transformations, and Symmetry}

We are continuing, not beginning, the discussion of Poisson brackets. In this lecture the subject is widened from a useful formula to a piece of the intrinsic structure of classical mechanics: phase space is a structured space, Poisson brackets describe flows on that space, and symmetries are to be understood as special flows that preserve the Hamiltonian. These notes follow Leonard Susskind's Lecture 9 closely, with the supporting frame evidence retained where the boardwork matters; the lecture is due to Leonard Susskind, and the supporting curation used here comes through LazyingArt LLC and Video2Book.

\section{Poisson Structure and the Idea of Flow}

Phase space is not just a pile of coordinates. It is not merely a list of \(q\)'s and \(p\)'s written next to one another. It is a space with structure, and what matters physically is the structure that survives the transformations we are allowed to perform.

A familiar comparison is metric structure. On a Riemannian space one has a notion of distance between neighboring points, and that metric data is part of what the space is. The Poisson structure of phase space is different. It is not about distance; it is about the special algebraic relations among the phase-space coordinates and about the transformations that preserve those relations.

This is why the lecture opens with a geometric widening of the frame. Before speaking about phase space itself, one may already think of ordinary rotations and translations in coordinate space as flows. A rotation can be built from many small rotations. A translation can be built from many small translations. One passes continuously through a family of nearby coordinate systems. In that sense, even a familiar symmetry in ordinary space is a flow.

Phase space inherits that idea and enlarges it. Now we are no longer restricted to motions that shuffle the coordinates \(x\), \(y\), and \(z\). We may consider transformations that mix positions and momenta. That is the richer setting in which Poisson brackets live.

One such flow is the actual time evolution of a system under a Hamiltonian. But it is not the only one. A symmetry can also be viewed as a flow, even when it is not the physical time development of the system. This is the first real sales pitch of the lecture: Poisson brackets are the language of flows on phase space.

With that motivation in place, the lecturer does not want to begin in the middle. So he deliberately backs up and reviews the subject abstractly before returning to concrete calculations.

\section{Abstract Rules for Poisson Brackets}

At this point we do what Susskind jokingly calls the French-mathematician thing: we abstract. Let \(A\), \(B\), and \(C\) be functions on phase space. Then the Poisson bracket is characterized first by general algebraic rules.

The first is antisymmetry:
\begin{equation}
\{A,B\}=-\{B,A\}.
\end{equation}

The second is linearity:
\begin{equation}
\{\alpha A+\beta B,C\}=\alpha\{A,C\}+\beta\{B,C\},
\end{equation}
where \(\alpha\) and \(\beta\) are ordinary numbers.

The third is the product rule:
\begin{equation}
\{AB,C\}=A\{B,C\}+B\{A,C\}.
\end{equation}

The lecture spends a little time on a point that is easy to rush past. These functions are ordinary numerical-valued functions on phase space. In classical mechanics their multiplication is commutative. So if an expression contains \(A\) times a Poisson bracket, the order of that multiplication does not matter. What does matter is the order inside the bracket itself. Reversing the bracket reverses the sign. That warning will matter again when the lecture turns Poisson brackets into differentiation rules.

One may allow complex-valued functions if one likes, but the essential point is simpler than that: outside the bracket we are multiplying commuting functions.

Only after this abstract review do we return to the concrete derivative formula. When needed, the bracket is given by
\begin{equation}
\{A,B\}=\sum_i\left(
\frac{\partial A}{\partial q_i}\frac{\partial B}{\partial p_i}
-
\frac{\partial A}{\partial p_i}\frac{\partial B}{\partial q_i}
\right).
\end{equation}
As in the lecture, we assume the functions are smooth enough for these manipulations to make sense.

\section{Concrete Canonical Brackets and Differentiation}

The abstract rules now have to meet the canonical variables themselves. At this stage of the lecture the board uses capital \(Q\)'s and \(P\)'s abstractly for canonical coordinates. Later, when explicit transformed variables are introduced, we will reserve capitals for the transformed coordinates and return to lower-case \(q\)'s and \(p\)'s for the original ones.

The surviving board frame catches the exact moment when the lecture moves from the general product rule to the self-bracket of a canonical momentum.

\begin{figure}[h]
\centering
\includegraphics[width=0.56\textwidth]{lecture_09_figure_02.png}
\caption{Poisson bracket of \(P\) with itself}
\end{figure}

The upper line in the frame is cropped, but the nearby context makes clear that it is the product rule
\begin{equation}
\{AB,C\}=A\{B,C\}+B\{A,C\},
\end{equation}
while the lower line isolates the concrete special case
\begin{equation}
\{P,P\}=0.
\end{equation}
This follows immediately from antisymmetry: a bracket equal to its own negative must vanish. The same argument gives
\begin{equation}
\{Q,Q\}=0.
\end{equation}

What antisymmetry does \emph{not} determine is the mixed bracket. For that one must use the derivative definition. In many degrees of freedom one finds
\begin{equation}
\{Q_i,P_j\}=\delta_{ij},\qquad \{P_i,Q_j\}=-\delta_{ij},
\end{equation}
together with
\begin{equation}
\{Q_i,Q_j\}=0,\qquad \{P_i,P_j\}=0.
\end{equation}

Two immediate consequences are worth recording. If both functions depend only on the coordinates, then their Poisson bracket vanishes; if both depend only on the momenta, it also vanishes:
\begin{equation}
\{F(Q),G(Q)\}=0,\qquad \{F(P),G(P)\}=0.
\end{equation}
To get something nonzero, a \(Q\) on one side must meet a \(P\) on the other.

\subsection{Question \& Answer}

Question. Why does antisymmetry immediately force \(\{P,P\}=0\), but not yet determine \(\{Q,P\}\)?

Answer. Because for a self-bracket we have
\[
\{P,P\}=-\{P,P\},
\]
so the bracket equals its own negative and must be zero. But for two different variables antisymmetry only gives
\[
\{Q,P\}=-\{P,Q\}.
\]
That is a relation between two different brackets, not a value. The value comes from the concrete derivative definition, and that is where the canonical normalization enters.

The lecture now turns that normalization into a useful operator identity. The second surviving frame shows the left-hand side of the next claim exactly where the lecturer writes it.

\begin{figure}[h]
\centering
\includegraphics[width=0.56\textwidth]{lecture_09_figure_03.png}
\caption{Poisson bracket with \(P\) as a derivative}
\end{figure}

Only the left-hand side and the equality sign are legible in the frame, but the transcript fixes the intended statement:
\begin{equation}
\{F(q),P\}=\frac{dF}{dq}.
\end{equation}
The lecture first states this, and only then proves it.

The proof is not done by a one-line appeal to the definition. Instead, it is done pedagogically, by working through monomials. By linearity, it is enough to prove the statement for monomials in \(q\).

First take \(F(q)=1\). Then the claim is \(\{1,P\}=0\). The product rule gives
\begin{equation}
\{1^2,P\}=1\{1,P\}+1\{1,P\}=2\{1,P\}.
\end{equation}
But \(1^2=1\), so the left-hand side is just \(\{1,P\}\). Therefore \(\{1,P\}=2\{1,P\}\), which forces
\begin{equation}
\{1,P\}=0.
\end{equation}

Next take \(F(q)=q\). Then
\begin{equation}
\{q,P\}=1,
\end{equation}
which is just the fundamental canonical bracket.

Now take \(F(q)=q^2\). The product rule gives
\begin{equation}
\{q^2,P\}=q\{q,P\}+\{q,P\}q=2q,
\end{equation}
which is exactly the derivative of \(q^2\). The induction to \(q^n\) is then straightforward. Linearity extends the result to polynomials, and from there the lecture passes to smooth functions by the usual approximation argument.

Thus Poisson bracketing with the conjugate momentum is another way of differentiating. In many-variable notation the two basic differentiation rules may be written
\begin{equation}
\{F(q_1,\dots,q_n),p_i\}=\frac{\partial F}{\partial q_i},
\qquad
\{q_i,F(p_1,\dots,p_n)\}=\frac{\partial F}{\partial p_i}.
\end{equation}
Here again the lecture insists on the order. To differentiate with respect to \(q_i\), the momentum goes on the right. To differentiate with respect to \(p_i\), the coordinate goes on the left.

\section{Time Evolution as a Poisson Flow}

Only now does the lecture add the dynamical postulate
\begin{equation}
\dot A=\{A,H\},
\end{equation}
for any function \(A\) on phase space.

The interpretation is geometric. A point moves along its orbit in phase space under the Hamiltonian flow. As it moves, the value of \(A(q,p)\) changes. That rate of change is the Poisson bracket of \(A\) with the Hamiltonian. In this way the time development of the system is itself just one special case of a flow generated by a function on phase space.

A student asks whether this already assumes that the \(P\)'s and \(Q\)'s are functions of time. The lecture's answer is practical: once we apply the rule to \(P\) and \(Q\) themselves, they become time-dependent by construction.

Take, for example,
\begin{equation}
H=\frac{P^2}{2M}.
\end{equation}
Then
\begin{equation}
\dot P=\left\{P,\frac{P^2}{2M}\right\}=0,
\end{equation}
because any function of \(P\) has vanishing Poisson bracket with any other function of \(P\). So momentum is conserved.

For the coordinate we get
\begin{align}
\dot Q
&=\left\{Q,\frac{P^2}{2M}\right\}
=\frac{1}{2M}\{Q,P^2\} \\
&=\frac{1}{2M}\left(P\{Q,P\}+\{Q,P\}P\right)
=\frac{1}{2M}(P+P)
=\frac{P}{M}.
\end{align}
So the familiar velocity-momentum relation reappears directly from the bracket algebra.

This is the bridge from structure to mechanics. Without importing Hamilton's equations from outside, we recover them from the Poisson algebra itself. The lecture also pauses here to say something stronger in tone, if not in theorem form: these basic formulas are not parochial. They are part of the structural core of classical mechanics.

\section{Canonical Transformations from Examples}

The next turn in the lecture is a wider one. If Hamiltonian time evolution is a flow on phase space, then what other changes of phase-space coordinates preserve the Poisson structure? This is not yet the question of symmetry for a given Hamiltonian. It is the more basic question: which transformations preserve the form of classical mechanics itself?

To keep signs under control, let us test the transformed variables in the normalized order
\[
\{Q,P\}=1.
\]
Now take a single canonical pair \((q,p)\) and try explicit examples.

First,
\begin{equation}
Q=2q,\qquad P=2p.
\end{equation}
Then
\begin{equation}
\{Q,P\}=\{2q,2p\}=4\{q,p\}=4,
\end{equation}
so the canonical normalization is not preserved. Stretching both directions by the same factor ruins the Poisson structure.

Try instead
\begin{equation}
Q=2q,\qquad P=\frac12 p.
\end{equation}
Now
\begin{equation}
\{Q,P\}=\left\{2q,\frac12 p\right\}=\{q,p\}=1.
\end{equation}
This works. One direction has been stretched while the other has been squeezed by the compensating factor. In the one-pair picture, the area is preserved.

The lecture pauses over that geometric intuition. For one \(p\) and one \(q\), area preservation is indeed the right picture. But once there are many pairs, the full condition is more subtle. Area preservation is the first intuition, not the whole story.

A second successful example is a rotation in phase space:
\begin{equation}
P=\cos\theta\,p+\sin\theta\,q,\qquad
Q=-\sin\theta\,p+\cos\theta\,q.
\end{equation}
By antisymmetry we already know
\[
\{P,P\}=0,\qquad \{Q,Q\}=0.
\]
So only the mixed bracket must be checked. Expanding by linearity,
\begin{align}
\{Q,P\}
&=\{-\sin\theta\,p+\cos\theta\,q,\cos\theta\,p+\sin\theta\,q\} \\
&=\cos^2\theta\,\{q,p\}-\sin^2\theta\,\{p,q\} \\
&=\cos^2\theta+\sin^2\theta \\
&=1.
\end{align}
So the phase-space rotation preserves the canonical bracket relations as well.

\begin{definition}
A canonical transformation is a change of phase-space coordinates
\[
(q_i,p_i)\longmapsto (Q_i,P_i)
\]
that preserves the canonical Poisson-bracket relations
\[
\{Q_i,Q_j\}=0,\qquad \{P_i,P_j\}=0,\qquad \{Q_i,P_j\}=\delta_{ij}.
\]
\end{definition}

This definition is algebraic, but the lecture adds an important conceptual restriction: we are talking about new coordinates for the same system. A canonical transformation is a re-description of the same phase space, not a mysterious map from one physical system into a completely different one. Nor is it a gauge redundancy. It is a change of phase-space coordinates that leaves the Poisson structure intact. If the Hamiltonian is rewritten in the new variables, the equations of motion retain their canonical form.

\subsection{Question \& Answer}

Question. What kinds of changes of \(p\) and \(q\) preserve mechanics, and why does area preservation enter?

Answer. In the one-pair case, the successful examples preserve area in phase space, and that is why the squeeze-stretch and the rotation work while the double stretch fails. But the true criterion is not geometric area alone; it is preservation of the canonical Poisson-bracket relations. Area preservation is the right first picture when there is one \(p\) and one \(q\), but the full many-variable condition is the canonical one.

\section{Infinitesimal Canonical Transformations and Generators}

At this point a student asks a natural question: is there any more useful way to express the condition for a canonical transformation than simply checking that all the Poisson brackets are preserved? The lecture answers yes, but cautiously. Not every global statement is made with full generality. The really useful construction is infinitesimal.

So let us build the transformation out of small steps:
\begin{equation}
Q=q+\delta q,\qquad P=p+\delta p.
\end{equation}
Here \(\delta q\) and \(\delta p\) are small functions of \(q\) and \(p\). Small means that quadratic and higher powers may be neglected.

Now demand that the transformed variables still satisfy the canonical relation. To first order,
\begin{align}
\{Q,P\}
&=\{q+\delta q,p+\delta p\} \\
&=\{q,p\}+\{\delta q,p\}+\{q,\delta p\}+\{\delta q,\delta p\}.
\end{align}
The last term is already quadratic in small quantities, so we drop it. Therefore the first-order canonical condition is
\begin{equation}
\{\delta q,p\}=-\{q,\delta p\}.
\end{equation}

This is still a condition to be checked. The lecture now turns it into a construction. Introduce a small parameter \(\epsilon\) and a function \(G(q,p)\), the generator, and define
\begin{equation}
\delta q=\epsilon\{q,G\},\qquad \delta p=\epsilon\{p,G\}.
\end{equation}
Using the differentiation rules just established, this becomes
\begin{equation}
\delta q=\epsilon\,\frac{\partial G}{\partial p},
\qquad
\delta p=-\epsilon\,\frac{\partial G}{\partial q}.
\end{equation}

Now the canonical condition becomes automatic. Indeed,
\[
\{\delta q,p\}
=\epsilon\left\{\frac{\partial G}{\partial p},p\right\}
=\epsilon\,\frac{\partial^2 G}{\partial q\,\partial p},
\]
while
\[
-\{q,\delta p\}
=-\left\{q,-\epsilon\frac{\partial G}{\partial q}\right\}
=\epsilon\,\frac{\partial^2 G}{\partial p\,\partial q}.
\]
These are the same mixed second derivative, so the first-order condition is satisfied identically.

What has happened conceptually is more important than the algebra. An infinitesimal canonical transformation is a flow on phase space. Every point is shifted a little bit. The function \(G\) tells us how that shift takes place.

Even better, the same generator controls the change of every phase-space function. If \(A(q,p)\) is any function, then
\begin{align}
\delta A
&=\frac{\partial A}{\partial q}\,\delta q+\frac{\partial A}{\partial p}\,\delta p \\
&=\epsilon\left(
\frac{\partial A}{\partial q}\frac{\partial G}{\partial p}
-
\frac{\partial A}{\partial p}\frac{\partial G}{\partial q}
\right) \\
&=\epsilon\{A,G\}.
\end{align}
So Poisson bracketing with \(G\) does not merely test the flow; it generates it.

This is the point at which Hamiltonian motion is reabsorbed into the larger picture. If we write \(\epsilon=\delta t\), then Hamiltonian evolution is simply the special case \(G=H\). In that sense the time evolution of the system is a canonical transformation generated by the Hamiltonian.

The lecture makes a strong-looking general remark here, then immediately softens it. The safe version is the one we keep: every infinitesimal canonical transformation is generated by some \(G\), and the finite canonical transformations of practical interest are built from such infinitesimal steps.

\subsection{Question \& Answer}

Question. If \(G\) generates a flow exactly the way \(H\) does, what distinguishes a general generator from the actual Hamiltonian?

Answer. Mathematically, very little. The same formalism generates both flows. A generator can be treated as though it were a Hamiltonian for the purpose of producing a motion on phase space. But the actual Hamiltonian of a physical system is not selected by mathematics alone. It is selected by the physics of the system, ultimately by experiment.

\section{Symmetry, Conservation, and the Free Particle}

Only now does the lecture come to symmetries themselves. If canonical transformations are the transformations that preserve the structure of mechanics, then a symmetry is the subclass that also leaves the Hamiltonian unchanged.

One may picture this geometrically. The Hamiltonian \(H\) generates the actual time flow of the system through phase space. Another generator \(G\) defines another flow, not necessarily the physical time flow. In general that \(G\)-flow will move the system from one energy surface to another. But if the \(G\)-flow happens to be tangent to the surfaces of constant energy, then it is a symmetry flow.

In formulas, the change of any function under the \(G\)-flow is
\[
\delta A=\epsilon\{A,G\}.
\]
Taking \(A=H\), the condition that the energy not change is
\begin{equation}
\delta H=\epsilon\{H,G\}=0,
\end{equation}
so the symmetry condition is
\begin{equation}
\{H,G\}=0.
\end{equation}
By antisymmetry this is equivalent to
\begin{equation}
\{G,H\}=0.
\end{equation}
But \(\{G,H\}\) is exactly the time derivative of \(G\), so
\begin{equation}
\dot G=\{G,H\}=0.
\end{equation}
Thus the generator of a symmetry is conserved under time evolution. The lecture dwells on the reciprocity: the same vanishing Poisson bracket says that \(G\) does not change \(H\), and also that \(H\) does not change \(G\).

\subsection{Question \& Answer}

Question. Why does ``the flow generated by \(G\) leaves \(H\) unchanged'' turn into ``\(G\) is conserved''?

Answer. Because it is the same equation read in the two possible orders. ``\(G\) leaves \(H\) unchanged'' means \(\{H,G\}=0\). Antisymmetry turns this into \(\{G,H\}=0\). But \(\{G,H\}\) is precisely \(\dot G\). The symmetry statement and the conservation statement are therefore two faces of the same vanishing Poisson bracket.

The lecture closes by making the abstract discussion concrete. In two dimensions take the free-particle Hamiltonian
\begin{equation}
H=\frac{p_x^2+p_y^2}{2m},
\end{equation}
and the angular momentum
\begin{equation}
G=xp_y-yp_x.
\end{equation}
Susskind temporarily sets \(m=1\) at the board, but the factor is inessential, so we keep it explicit.

Now compute:
\begin{align}
\{G,H\}
&=\frac{1}{2m}\{xp_y-yp_x,p_x^2+p_y^2\} \\
&=\frac{1}{2m}\left(
\{xp_y,p_x^2\}
+\{xp_y,p_y^2\}
-\{yp_x,p_x^2\}
-\{yp_x,p_y^2\}
\right).
\end{align}
At this stage the lecture proceeds term by term, using only the basic brackets
\[
\{x,p_x\}=1,\qquad \{y,p_y\}=1,
\]
with all cross-brackets vanishing.

The mixed terms \(\{xp_y,p_y^2\}\) and \(\{yp_x,p_x^2\}\) vanish immediately, because no \(x\) meets a \(p_x\) in the first and no \(y\) meets a \(p_y\) in the second. For the remaining two terms, the product rule gives
\begin{align}
\{xp_y,p_x^2\}
&=p_y\{x,p_x^2\}+x\{p_y,p_x^2\}
=2p_xp_y, \\
\{yp_x,p_y^2\}
&=p_x\{y,p_y^2\}+y\{p_x,p_y^2\}
=2p_xp_y.
\end{align}
Therefore
\begin{equation}
\{G,H\}
=\frac{1}{2m}\left(2p_xp_y-2p_xp_y\right)=0.
\end{equation}
So angular momentum is conserved for the free particle.

The lecture immediately turns the reciprocity around. Suppose we invent a new system whose Hamiltonian is
\begin{equation}
H=xp_y-yp_x.
\end{equation}
Then Hamilton's equations give
\begin{equation}
\dot x=\{x,H\}=-y,\qquad \dot y=\{y,H\}=x.
\end{equation}
That is circular motion. And because the Poisson bracket with \(p_x^2+p_y^2\) vanishes, that quantity is conserved in this new system. Whichever member of a vanishing Poisson-bracket pair is chosen as the Hamiltonian, the other becomes a conserved quantity.

One final remark belongs to the tone of the lecture. At several moments Susskind wants to say ``commute,'' then stops himself. In quantum mechanics that language will be irresistible. Here we keep the statement classical: the Poisson bracket vanishes.

\section{Summary}

The chapter has unfolded in the same rhythm as the lecture. We began by widening the subject from a bracket formula to the structure of phase space itself. We then reviewed the abstract algebra of Poisson brackets, descended to the canonical relations and the differentiation identities, added the time-evolution rule \(\dot A=\{A,H\}\), and finally widened again to canonical transformations and their generators.

The reward for keeping that order is that symmetry and conservation emerge naturally from the flow picture. A generator \(G\) defines a canonical flow. If that flow leaves the Hamiltonian unchanged, then \(\{H,G\}=0\), and the same equation says that \(G\) is conserved under Hamiltonian time evolution. This is highly abstract in classical mechanics, and the lecture ends by hinting that the same structure becomes even more concrete in quantum mechanics. But even here, in classical mechanics, the point is already visible: the Poisson structure of phase space is what lets mechanics recognize itself under change.